% Scientific Reports / Springer Nature layout
% Manuscript structure intentionally retained from the original paper:
% Introduction -> Related Work -> Method -> Experiments -> Limitations and Future Work -> Conclusion
\documentclass[pdflatex,sn-nature]{sn-jnl}

\usepackage{graphicx}
\usepackage{multirow}
\usepackage{amsmath,amssymb,amsfonts}
\usepackage{booktabs}
\usepackage{adjustbox}
\usepackage{placeins}
\usepackage{siunitx}
\usepackage{xcolor}
\newcommand{\rev}[1]{\textcolor{red}{#1}}

\newcommand{\TP}{\mathrm{TP}}
\newcommand{\FP}{\mathrm{FP}}
\newcommand{\FN}{\mathrm{FN}}
\newcommand{\Precision}{\mathrm{Precision}}
\newcommand{\Recall}{\mathrm{Recall}}

% Keep the official template's flexible page bottom. Top-preferred floats reduce
% large blank regions while allowing figures/tables to remain close to first mention.
\raggedbottom

\begin{document}

\title[FFRM-YOLO]{FFRM-YOLO: Frequency-Decoupled Feature Refinement for Metal Surface Defect Detection}

% Replace the placeholders below with the actual author/affiliation information.
\author*[1,2]{\fnm{First} \sur{Author}}\email{iauthor@gmail.com}

\author[2,3]{\fnm{Zeng} \sur{Jie}}\email{13797343752@163.com}
\equalcont{These authors contributed equally to this work.}

\author[2]{\fnm{Tao} \sur{Jian}}\email{taojian@zjut.edu.cn}

\affil*[1]{\orgdiv{Department}, \orgname{Organization}, \orgaddress{\street{Street}, \city{City}, \postcode{100190}, \state{State}, \country{Country}}}

\affil[2]{\orgdiv{College of Computer Science and Technology}, \orgname{Zhejiang University of Technology}, \orgaddress{\city{Hangzhou}, \postcode{310023}, \state{Zhejiang}, \country{China}}}

\affil[3]{\orgdiv{Department}, \orgname{Organization}, \orgaddress{\street{Street}, \city{City}, \postcode{610101}, \state{State}, \country{Country}}}

\abstract{Metal surface defect detection requires effective representations of both structural variations and fine texture details, yet spatial-domain methods often lack explicit modeling of their distinct frequency characteristics. To address this limitation, we propose FFRM-YOLO, a frequency-domain feature refinement framework built upon YOLO11n. The framework incorporates a frequency-domain feature refinement module (FFRM) that uses the discrete wavelet transform to separate features into low-frequency structural components and high-frequency directional details for specialized enhancement. In the low-frequency branch, a global–local recalibration module (GLRM) combines Mamba-SS2D with local depthwise convolutions to capture long-range dependencies while preserving local structural information. In the high-frequency branch, a multi-scale directional refinement module (MDRM) enhances directional textures through multi-scale modeling and cross-subband interaction. Additionally, Inner-CIoU loss is adopted to improve bounding-box regression.\rev{Experiments on three metal surface defect datasets show that FFRM-YOLO achieves mAP50 scores of \(78.6 \pm 0.5\%\), \(66.1 \pm 0.3\%\), and \(60.6 \pm 0.6\%\) on NEU-DET, GC10-DET, and APDDD, respectively, over three independent training runs. These results exceed the corresponding mean values of YOLO11n by 2.9, 2.7, and 3.0 percentage points, respectively.}The model requires only 4.8 M parameters and 7.7 GFLOPs. These results demonstrate the effectiveness of differentiated frequency modeling across diverse defect morphologies and data distributions, with a favorable balance between detection accuracy and computational efficiency.}

\keywords{Metal surface defect detection, Frequency-domain feature refinement, Discrete wavelet transform, State space model, Global--local feature modeling, Directional feature refinement.}

\maketitle

\section{Introduction}\label{sec:introduction}

    Metal surface defects arising during manufacturing can compromise product quality and subsequent processing. Despite advances in deep-learning-based detection, substantial morphological variations, blurred boundaries, and interference from background textures continue to pose challenges. Accurate detection requires representations that capture both the overall structure of defect regions and fine local details, whose relative importance varies with defect morphology. Effectively preserving and integrating these complementary cues is therefore an important challenge in metal surface defect detection.
    
    In recent years, single-stage object detectors, exemplified by YOLO, have been widely employed in industrial defect detection owing to their end-to-end training paradigm and high inference efficiency \cite{Redmon2016YOLO,Jocher2024YOLO11}. Existing studies have improved detection performance through multi-scale feature fusion, attention mechanisms, detection head refinement, and lightweight network design \cite{Yu2022ESNet,Liu2025GCNet,Liu2026GSSAYOLO,Zheng2026ACMSPYOLO,Zhang2026DRFMFF,Zhang2024MSFCRA}. However, these improvements primarily focus on spatial feature processing, without explicitly separating frequency components for subband-specific enhancement. From a frequency perspective, broad intensity variations within defect regions are generally associated with low-frequency components, whereas abrupt transitions and fine textures are more prominent in high-frequency components. \rev{Low-frequency components primarily characterize broad structural variations, whereas high-frequency components emphasize localized transitions, boundaries, and fine textures. These complementary characteristics may coexist within the same defect and vary across defect morphologies, rather than establishing a strict correspondence between individual defect categories and specific frequency bands.} They motivate explicit frequency decomposition and subband-specific enhancement of structural information and texture details.

    
    The discrete wavelet transform (DWT) decomposes features into a low-frequency approximation subband and three high-frequency detail subbands, providing spatially localized representations of frequency information \cite{Mallat1989Wavelet,Liu2018MWCNN}. Several studies have incorporated the wavelet transform into visual recognition and industrial defect analysis to replace downsampling, construct frequency-domain attention, or enhance edge information \cite{Cui2025SWT,Chen2026WTCFNet}. Meanwhile, state space models provide an efficient alternative for long-range dependency modeling with linear computational complexity \cite{Gu2024Mamba,Liu2024VMamba}. Recent approaches have further attempted to jointly utilize wavelet features and state space models for anomaly detection and industrial defect detection \cite{Zhang2025WaveMambaAD,Chen2026MWYOLO,Ge2026WaveMambaYOLO}. These developments motivate a more specific design question: how can global structural modeling and directional detail enhancement be coordinated through subband-specific processing within a wavelet decomposition and reconstruction framework? In particular, an effective design should capture long-range dependencies in low-frequency features while exploiting the directional priors of the high-frequency subbands and allowing information exchange among them.
    
    To address the above issues, this work proposes FFRM-YOLO, a metal surface defect detector with frequency-domain feature refinement. Taking YOLO11n as the baseline, the proposed model replaces the C3k2 modules in the backbone and neck with the frequency-domain feature refinement module (FFRM). FFRM uses the Haar DWT to separate features into low-frequency and high-frequency subbands and processes them through two specialized branches before inverse wavelet reconstruction. In the low-frequency branch, a global--local recalibration module (GLRM) combines the 2D selective scan (SS2D) module introduced in VMamba with local depthwise convolutions to capture long-range dependencies and preserve local structural information. Applying SS2D to the LL subband reduces its computational cost compared with applying it to the original-resolution feature map. In the high-frequency branch, a multi-scale directional refinement module (MDRM) combines direction-aware multi-scale processing, cross-subband interaction, and spatial gating to enhance texture and edge features by exploiting the directional priors of different subbands. The refined subbands are reconstructed through the inverse wavelet transform and integrated with the input through a residual connection. Inner-CIoU~\cite{Zhang2023InnerIoU} is additionally adopted as the bounding-box regression loss.\rev{Experiments on NEU-DET, GC10-DET, and APDDD demonstrate improvements in mean detection performance over YOLO11n across three independent training runs, with 4.8 M parameters and 7.7 GFLOPs.}
    
    The main contributions of this work are as follows:
    
    \begin{enumerate}
    \item We propose FFRM-YOLO, a frequency-aware metal surface defect detector that integrates subband-specific enhancement of low-frequency and high-frequency features into YOLO11n within a wavelet decomposition and reconstruction framework.
    
    \item We design GLRM for the low-frequency subband, combining SS2D-based global modeling with local depthwise convolutions to enhance structural representations at reduced spatial resolution.
    
    \item We develop MDRM for the high-frequency subbands, integrating direction-aware multi-scale processing and cross-subband interaction to enhance texture and edge features while exploiting the directional priors of different subbands.
    \end{enumerate}

\section{Related Work}\label{sec:related}
\subsection{Industrial Defect Detection}
Deep-learning-based industrial defect detection methods can be mainly categorized into single-stage, two-stage, and Transformer-based end-to-end approaches. Single-stage methods, exemplified by the YOLO series, have been widely applied to industrial inspection owing to their high efficiency in training and inference; for instance, GSSA-YOLO \cite{Liu2026GSSAYOLO} and ACMSP-YOLO \cite{Zheng2026ACMSPYOLO} achieve parameter compression and accuracy improvements in the detection of automotive steel and steel strips through lightweight spatial kernel attention and shared multi-semantic spatial attention, respectively. Two-stage methods, represented by Cascade R-CNN \cite{Cai2018CascadeRCNN} and Mask R-CNN \cite{He2017MaskRCNN}, extend the applicability of the two-stage framework in terms of cascaded regression and instance segmentation. Transformer-based end-to-end detectors, represented by DETR \cite{Carion2020DETR}, have advanced rapidly in recent years: SH-DETR \cite{Wu2025SHDETR} attains high accuracy on steel defect datasets through heterogeneous multi-scale feature enhancement, while MESC-DETR \cite{Zhou2025MESCDETR}, RT-LSM-DETR \cite{Hsiang2025RTLSMDETR}, HMA-DETR \cite{Yu2026HMADETR}, and Lite-DETR \cite{Luo2025LiteDETR} improve DETR from the perspectives of multi-scale defect adaptation, lightweight state space modeling, hybrid adaptive features, and lightweight architectures, respectively. Although the aforementioned methods enhance detection performance under different architectural paradigms, their feature extraction is still conducted primarily in the spatial domain, where features with distinct frequency responses are typically processed with identical or similar operations. Therefore, this study introduces the discrete wavelet transform (DWT) to establish a frequency-domain decoupling perspective and designs the FFRM to route heterogeneous frequency subbands into differentiated processing pathways, thereby overcoming the inherent limitations of existing methods in feature extraction.

\subsection{Wavelet Transform in Visual Recognition}
Owing to its lightweight and invertible signal-analysis characteristics, the discrete wavelet transform (DWT) has become an important tool for frequency-domain enhancement in deep networks. In early work, MWCNN \cite{Liu2018MWCNN} embeds multi-level wavelet decomposition into convolutional networks, validating the effectiveness of wavelet-domain modeling for image restoration. Subsequently, SWT-based networks \cite{Cui2025SWT} and EWT \cite{Li2024EWT} exploit the stationary wavelet transform and wavelet-domain sparse representations to improve defect recognition and image denoising performance, respectively. In low-level vision, Wave-Mamba \cite{Zou2024WaveMamba} replaces strided downsampling with the DWT, constructing a lossless frequency-decoupling operator. In the field of industrial inspection, WE-PaDiM \cite{Gardner2026WEPaDiM} demonstrates that DWT subband decomposition can facilitate feature selection for anomaly detection, while Wave-MambaAD \cite{Zhang2025WaveMambaAD} couples wavelet decomposition with state space scanning to achieve more competitive anomaly detection performance at a lower cost. Beyond anomaly detection, WTCFNet \cite{Chen2026WTCFNet} introduces wavelet features for industrial defect detection, whereas WFSCA-YOLO \cite{Yan2026WFSCAYOLO} explores frequency--spatial interaction for optical object detection under atmospheric degradation. More recently, WaveMamba-YOLO \cite{Ge2026WaveMambaYOLO} combines frequency-aware decomposition with state space modeling for surface defect localization. However, existing methods still exhibit limitations in fine-grained subband modeling. Although some approaches explicitly separate low- and high-frequency components, the heterogeneous characteristics of different wavelet subbands are not fully exploited. In particular, the distinct scale responses and complementary relationships among horizontal, vertical, and diagonal high-frequency components remain insufficiently modeled. To address these limitations, this work adopts a dual-branch architecture within the FFRM, where the GLRM focuses on global low-frequency structural modeling, while the MDRM performs multi-scale refinement and cross-subband interaction for high-frequency directional details.


\subsection{Efficient Visual Representation With State Space Models}
Owing to their linear-time complexity and capability of capturing long-range dependencies, state space models (SSMs) are emerging as efficient alternatives to Transformers. Mamba \cite{Gu2024Mamba} achieves a breakthrough in natural language processing (NLP) tasks through selective state space modeling. VMamba \cite{Liu2024VMamba} introduces the two-dimensional selective scan (SS2D) and constructs a general-purpose visual backbone with linear complexity. LocalMamba \cite{Huang2024LocalMamba} further reveals that global flattened scanning disrupts the local structure of visual features and proposes a windowed scanning strategy. In industrial inspection, MambaAD \cite{He2024MambaAD} realizes multi-class unsupervised anomaly detection via a local--global decoder, while YOLO-GVMamba \cite{Lei2025YOLOGVMamba} explores real-time steel surface defect detection and MobileMamba \cite{He2025MobileMamba} investigates efficient SSM integration for lightweight visual recognition. Recent studies have further extended lightweight state-space modeling to different detection scenarios. SCS-Mamba \cite{Xu2026SCSMamba} focuses on underwater object detection, whereas LVMSOD \cite{Fu2026LVMSOD} investigates small-object perception for autonomous vehicles. In industrial inspection, MWYOLO \cite{Chen2026MWYOLO} integrates Mamba with multi-frequency attention for lightweight surface defect detection. However, existing paradigms typically apply spatial scanning directly to full-resolution feature maps, leaving the potential of computational efficiency insufficiently exploited and lacking targeted designs for features at different frequencies. To this end, this work proposes the GLRM, which takes the half-resolution LL subband derived by the DWT as the scanning carrier, thereby achieving efficient global modeling of defect structural features while improving computational efficiency.


\section{Methodology}\label{sec:method}

\subsection{Overall Architecture}
As illustrated in Fig.~\ref{fig:architecture}, FFRM-YOLO builds upon YOLO11n \cite{Jocher2024YOLO11} by replacing all C3k2 modules in the backbone and neck with the proposed frequency-decoupled feature refinement module (FFRM). FFRM performs wavelet decomposition, subband-specific enhancement, and reconstruction to integrate structural information and texture details. Inner-CIoU is additionally adopted as the bounding-box regression loss during training.

\begin{figure}[!t]
        \centering
        \includegraphics[width=0.98\textwidth]{architecture.png}
            \caption{Overall architecture of FFRM-YOLO. B1--B4 denote the replaceable blocks in the backbone from shallow to deep. N1--N4 denote the replaceable blocks along the neck feature-fusion path, as labeled.}
            \label{fig:architecture}
    \end{figure}

Given an input feature map $\mathbf{X}\in\mathbb{R}^{B\times C\times H\times W}$, where $B$, $C$, $H$, and $W$ denote the batch size, channel number, height, and width, respectively, FFRM first applies the two-dimensional Haar DWT \cite{Mallat1989Wavelet}:
\begin{equation}
\left(
\mathbf{X}_{LL},
\mathbf{X}_{LH},
\mathbf{X}_{HL},
\mathbf{X}_{HH}
\right)
=
\mathrm{DWT}\!\left(\mathcal{P}(\mathbf{X})\right),
\label{eq:overview_dwt}
\end{equation}
where $\mathcal{P}(\cdot)$ applies zero padding to the bottom and right boundaries when necessary to obtain even spatial dimensions $H'$ and $W'$. Each subband has dimensions $B\times C\times H'/2\times W'/2$. The LL subband contains low-frequency approximation information, while LH, HL, and HH capture high-frequency details in different directions.

\begin{figure}[!t]
    \centering
    \includegraphics[width=0.98\textwidth]{FFRM.png}
    \caption{Architecture of the frequency-domain feature refinement module (FFRM). DWT decomposes input features into low- and high-frequency subbands, which are refined by GLRM and MDRM, respectively, followed by inverse wavelet reconstruction, channel projection, and residual fusion.}
    \label{fig:ffrm}
\end{figure}

As shown in Fig.~\ref{fig:ffrm}, the subbands are enhanced through two branches:
\begin{equation}
\begin{aligned}
\mathbf{Y}_{LL}
&= \mathrm{GLRM}(\mathbf{X}_{LL}), \\
\left[
\mathbf{Y}''_{LH},
\mathbf{Y}''_{HL},
\mathbf{Y}''_{HH}
\right]
&= \mathrm{MDRM}\!\left(
\left[
\mathbf{X}_{LH},
\mathbf{X}_{HL},
\mathbf{X}_{HH}
\right]\right),
\end{aligned}
\label{eq:overview_branches}
\end{equation}
where $[\cdot]$ denotes channel-wise concatenation. $\mathrm{GLRM}(\cdot)$ combines SS2D-based long-range modeling with local depthwise convolutions to enhance low-frequency structural information. $\mathrm{MDRM}(\cdot)$ employs direction-aware multi-scale processing, cross-subband interaction, and spatial attention to refine high-frequency texture and edge features. Its output is partitioned into three refined subbands in the same order as the input.

The refined subbands are reconstructed using the IDWT, followed by channel projection and residual fusion:
\begin{equation}
\begin{aligned}
\hat{\mathbf{Y}}
&= \mathcal{C}_{H,W}\!\left(
\mathrm{IDWT}\!\left(
\mathbf{Y}_{LL},
\mathbf{Y}''_{LH},
\mathbf{Y}''_{HL},
\mathbf{Y}''_{HH}
\right)\right), \\
\mathbf{Y}_{\mathrm{FFRM}}
&= \mathbf{X}
+ \mathrm{BN}\!\left(
\mathrm{Conv}_{1\times1}(\hat{\mathbf{Y}})
\right),
\end{aligned}
\label{eq:overview_reconstruction}
\end{equation}
where $\mathcal{C}_{H,W}(\cdot)$ removes any padding by cropping the reconstructed features to the original spatial size, and $\mathrm{BN}(\cdot)$ denotes batch normalization. The $1\times1$ convolution preserves the channel dimension, yielding $\mathbf{Y}_{\mathrm{FFRM}} \in\mathbb{R}^{B\times C\times H\times W}$.\rev{In implementation, the Haar DWT applies four fixed
$2\times2$ filters independently to each input channel
with stride~2, while the IDWT performs the corresponding
inverse Haar reconstruction. Padding and cropping are
applied only to the transformation path, ensuring that the
reconstructed features match the original input dimensions
before residual addition.}


\subsection{Global--Local Recalibration Module}
The GLRM enhances low-frequency structural features through parallel global and local branches, as illustrated in Fig.~\ref{fig:glrm}. The global branch uses 2D selective scan (SS2D) to capture long-range dependencies, while the local branch employs a depthwise convolution to capture local spatial variations. This design combines global context with local structural information at a low computational cost.

\begin{figure}[!t]
    \centering
            \includegraphics[width=0.98\textwidth]{GLRM.png}
        \caption{Architecture of the global--local recalibration module (GLRM). Parallel SS2D-based global modeling and local depthwise convolution refine low-frequency features, with their outputs fused through learnable scalar weights.}
        \label{fig:glrm}
    \end{figure}

\noindent\textbf{Global branch.} Given the low-frequency input $\mathbf{X}_{LL}$, layer normalization is first applied to obtain $\tilde{\mathbf{X}}_{LL}=\mathrm{LN}(\mathbf{X}_{LL})$. The global branch processes the normalized features using four directional scans: row-wise and column-wise traversals and their reverse orders. The resulting sequences are mapped back to their spatial arrangement and aggregated:

\begin{equation}
    \mathbf{Z}_{\mathrm{ssm}}
    =
    \frac{1}{4}\sum_{k=1}^{4}
    \mathcal{S}_{k}^{-1}\!\left(
    \mathcal{M}\!\left(
    \mathcal{S}_{k}(\tilde{\mathbf{X}}_{LL})
    \right)\right),
    \label{eq:glrm_global}
    \end{equation}
where $\mathcal{S}_{k}(\cdot)$ denotes the $k$-th scan operation, $\mathcal{S}_{k}^{-1}(\cdot)$ restores the corresponding spatial arrangement, and $\mathcal{M}(\cdot)$ denotes selective state space sequence processing. For a fixed channel dimension, this operation has linear complexity with respect to the number of spatial positions $N$ in the LL subband~\cite{Liu2024VMamba}.


\noindent\textbf{Local branch.} In parallel, the local branch applies a $7\times7$ depthwise convolution, followed by batch normalization and SiLU activation:
\begin{equation}
\mathbf{Z}_{\mathrm{loc}}
=
\mathrm{SiLU}\!\left(
\mathrm{BN}\!\left(
\mathrm{DWConv}_{7\times7}
(\tilde{\mathbf{X}}_{LL})
\right)\right).
\label{eq:glrm_local}
\end{equation}
This branch complements global sequence modeling by capturing local spatial patterns in the low-frequency features.

The two branch outputs are combined using learnable scalar weights, followed by a $1\times1$ convolution, batch normalization, and a residual connection:
\begin{equation}
\begin{aligned}
\mathbf{Z}_{\mathrm{fused}}
&=
w_{\mathrm{mamba}}\mathbf{Z}_{\mathrm{ssm}}
+
w_{\mathrm{local}}\mathbf{Z}_{\mathrm{loc}}, \\
\mathbf{Y}_{LL}
&=
\mathbf{X}_{LL}
+
\mathrm{SiLU}\!\left(
\mathrm{BN}\!\left(
\mathrm{PWConv}_{1\times1}
(\mathbf{Z}_{\mathrm{fused}})
\right)
\right).
\end{aligned}
\label{eq:glrm_fusion}
\end{equation}
Here, $w_{\mathrm{mamba}}$ and $w_{\mathrm{local}}$ are unconstrained learnable scalars initialized to $1$ and optimized jointly with the network parameters. The resulting features combine long-range context and local structural information for subsequent wavelet reconstruction.


\subsection{Multi-Scale Directional Refinement Module}

The MDRM refines the three high-frequency subbands through multi-scale spatial processing, channel recalibration, and spatial attention. As illustrated in Fig.~\ref{fig:mdrm}, it consists of multi-kernel depthwise convolution (MK-DeConv), cross-subband excitation (CSE), and high-frequency spatial attention (HFSA). These stages combine subband-specific feature extraction with information exchange and spatial refinement.

\begin{figure}[!t]
    \centering
    \includegraphics[width=0.98\textwidth]{MDRM.png}
    \caption{Architecture of the multi-scale directional refinement module (MDRM). High-frequency features are refined through subband-specific multi-scale processing, joint channel recalibration, and spatial attention. HFSA shares parameters across the three feature groups while generating a separate attention map for each group.}
    \label{fig:mdrm}
\end{figure}


\noindent\textbf{Multi-kernel depthwise convolution.}
The three input subbands $\mathbf{X}_{LH}$, $\mathbf{X}_{HL}$, and $\mathbf{X}_{HH}$ are processed separately using depthwise convolutions with kernel sizes of $3\times3$, $5\times5$, and $7\times7$, respectively:
\begin{equation}
\begin{aligned}
Y_{\mathrm{MK}}
&=
\mathrm{BN}\!\left(
\mathrm{PWConv}\!\left(
\left[
\psi_{3}(\mathbf{X}_{\mathrm{LH}}),\,
\psi_{5}(\mathbf{X}_{\mathrm{HL}}),\,
\psi_{7}(\mathbf{X}_{\mathrm{HH}})
\right]
\right)
\right), \\
\psi_k(\mathbf{X})
&=
\mathrm{SiLU}\!\left(
\mathrm{BN}\!\left(\mathrm{DWConv}_{k\times k}(\mathbf{X})\right)
\right).
\end{aligned}
\label{eq:mdrm_mk}
\end{equation}
Here, $[\cdot]$ denotes channel-wise concatenation. The different kernel sizes provide different receptive fields across the directional subbands. The depthwise convolutions $\mathrm{DWConv}_{k\times k}(\cdot)$ process each subband independently, while the subsequent $1\times1$ convolution $\mathrm{PWConv}(\cdot)$ mixes information across the concatenated channels and preserves the total channel dimension of $3C$.\rev{The directional characteristics originate from the LH, HL, and HH wavelet subbands, while the square depthwise convolutions provide subband-specific receptive fields rather than explicitly orientation-selective filtering. The subsequent pointwise convolution and channel recalibration facilitate information exchange among the refined high-frequency features.}

\noindent\textbf{Cross-subband excitation.}
Following the squeeze-and-excitation mechanism~\cite{Hu2018SENet},the CSE uses global channel descriptors to recalibrate the combined high-frequency features:
\begin{equation}
\begin{aligned}
\mathbf{Y}_{\mathrm{CSE}}
&=
\sigma\!\left(
\mathbf{W}_2
\mathrm{ReLU}\!\left(
\mathbf{W}_1
\mathrm{GAP}(
    \mathbf{Y}_{\mathrm{MK}})
\right)\right)\odot
\mathbf{Y}_{\mathrm{MK}},
\end{aligned}
\label{eq:mdrm_cse}
\end{equation}
where $\mathrm{GAP}(\cdot)$ denotes global average pooling, $\sigma(\cdot)$ is the sigmoid function, and $\odot$ denotes element-wise multiplication with broadcasting. The projection matrices are $\mathbf{W}_1\in\mathbb{R}^{d\times3C}$ and $\mathbf{W}_2\in\mathbb{R}^{3C\times d}$, where $d=3C/8$ corresponds to a reduction ratio of $8$. By computing channel weights jointly from the combined representation, CSE models channel dependencies across features originating from different subbands.


\noindent\textbf{High-frequency spatial attention.}
The recalibrated features are split into three $C$-channel
groups:
\begin{equation}
\left(
\mathbf{Y}_{LH}',
\mathbf{Y}_{HL}',
\mathbf{Y}_{HH}'
\right)
=
\mathrm{Split}\!\left(\mathbf{Y}_{\mathrm{CSE}}\right).
\label{eq:mdrm_split}
\end{equation}
These groups correspond to the three high-frequency inputs of the subsequent IDWT, with each group containing information from the preceding channel interaction.

First, the HFSA applies a shared spatial gating network to each group. For an input group $\mathbf{Y}_k'$, a $1\times1$ convolution first reduces the channel dimension from $C$ to $m$, where $m=\max(4,\min(8,\lfloor C/4\rfloor))$. A $3\times3$ convolution then captures local spatial information while maintaining $m$ channels, producing the intermediate feature $\mathbf{U}_k$. Both convolutions preserve the spatial dimensions.

Then, the intermediate feature is passed through GELU activation and a $1\times1$ convolution that reduces the channel dimension to one. A sigmoid function then produces the spatial attention map $\mathbf{A}_k$:
\begin{equation}
\mathbf{A}_k
=
\sigma\!\left(
\mathrm{Conv}_{1\times1}^{m\rightarrow1}\!\left(
\mathrm{GELU}(\mathbf{U}_k)
\right)\right),
\end{equation}
where $k\in\{LH,HL,HH\}$. Then, the attention map is broadcast across the $C$ channels and applied to the corresponding input group:
\begin{equation}
\mathbf{Y}_k''
=
\mathbf{A}_k\odot\mathbf{Y}_k',
\label{eq:mdrm_hfsa}
\end{equation}
where $\odot$ denotes element-wise multiplication. Note that HFSA shares its convolutional parameters across the three groups, but generates a separate attention map from each group's features. The refined outputs are subsequently combined with the low-frequency output through the IDWT.

\subsection{Training Losses}
\label{sec:training_losses}

FFRM-YOLO retains the binary cross-entropy (BCE) loss and distribution focal loss (DFL) of YOLO11n~\cite{Jocher2024YOLO11} for classification and bounding-box distribution learning, respectively. The original CIoU term is replaced with Inner-CIoU based on the auxiliary-box strategy of Inner-IoU~\cite{Zhang2023InnerIoU}.

\noindent\textbf{Classification loss.}
Binary cross-entropy (BCE) supervises the predicted class scores:
\begin{equation}
\mathcal{L}_{\mathrm{BCE}}(p,y)
=
-y\log p-(1-y)\log(1-p),
\label{eq:classification_loss}
\end{equation}
where $p$ denotes a predicted class probability and $y$ is its corresponding target score. The target assignment, weighting, and normalization follow the baseline implementation.

\noindent\textbf{Inner-CIoU loss.}
Inner-CIoU combines auxiliary-box overlap with the center-distance and aspect-ratio penalties of CIoU \cite{Zheng2020distance,Zhang2023InnerIoU}. For the predicted box $B_p$ and ground-truth box $B_g$, auxiliary boxes $B_p^\beta$ and $B_g^\beta$ are obtained by scaling their widths and heights by $\beta$ while keeping their centers unchanged. We set $\beta=0.70$ in all experiments. Their overlap is computed as:
\begin{equation}
\mathrm{IoU}_{\mathrm{inner}}
=
\frac{|B_p^\beta\cap B_g^\beta|}
{|B_p^\beta\cup B_g^\beta|},
\label{eq:inner_iou}
\end{equation}
where $|\cdot|$ denotes area. The regression loss is computed as:
\begin{equation}
\mathcal{L}_{\mathrm{Inner\text{-}CIoU}}
=
1-\mathrm{IoU}_{\mathrm{inner}}
+
\frac{\rho^2(\mathbf{c}_p,\mathbf{c}_g)}{d^2}
+
\alpha v,
\label{eq:inner_ciou}
\end{equation}
where $\rho$ is the Euclidean distance between the original box centers $\mathbf{c}_p$ and $\mathbf{c}_g$, and $d$ is the diagonal length of their smallest enclosing rectangle. The aspect-ratio term $v$ and its weight $\alpha$ follow the standard CIoU formulation and are computed from the original boxes.

\noindent\textbf{Distribution focal loss.}
DFL~\cite{Li2020GFL} supervises the discrete distributions of the distances from each reference point to the four bounding-box sides. For a target distance $t$ expressed in feature-map units, let $i=\lfloor t\rfloor$ and $i+1$ denote its adjacent discrete bins. The loss for one side is:
\begin{equation}
\mathcal{L}_{\mathrm{DFL}}(t)
=
-(i+1-t)\log q_i
-(t-i)\log q_{i+1},
\label{eq:dfl}
\end{equation}
where $q_i$ and $q_{i+1}$ are the predicted probabilities of the two bins. The loss is averaged over the four sides and aggregated over positive samples using the weighting and normalization of YOLO11n. DFL supervises the boundary-distance distributions, while Inner-CIoU measures the geometric agreement between decoded boxes and ground-truth boxes.

\noindent\textbf{Overall loss.}
The overall training objective is:
\begin{equation}
\begin{aligned}
\mathcal{L}_{\mathrm{total}}
={}&
\lambda_{\mathrm{cls}}\mathcal{L}_{\mathrm{BCE}}
+
\lambda_{\mathrm{box}}\mathcal{L}_{\mathrm{Inner\text{-}CIoU}} \\
&+
\lambda_{\mathrm{dfl}}\mathcal{L}_{\mathrm{DFL}},
\end{aligned}
\label{eq:total_loss}
\end{equation}
where the three terms denote the aggregated classification, box-overlap, and distribution losses, respectively. The loss weights follow the YOLO11n baseline configuration.

\section{Experiments}\label{sec:experiments}

This section describes the datasets, evaluation metrics, and implementation details, followed by comparative experiments, qualitative visualizations, and ablation studies.

\subsection{Datasets}
We evaluate FFRM-YOLO on three metal surface defect datasets: NEU-DET, GC10-DET, and APDDD. These datasets cover different materials, defect categories, and visual characteristics, enabling evaluation under diverse surface inspection conditions. Each dataset is randomly divided into training and test sets at a ratio of 7:3. The same split of each dataset is used for all methods.

\textbf{NEU-DET:} NEU-DET~\cite{Song2013NEU} is a benchmark for hot-rolled steel strip surface defect detection. It contains 1,800 grayscale images from six categories: Crazing, Inclusion, Patches, Pitted Surface, Rolled-in Scale, and Scratches, with 300 images per category. Intra-class variations, inter-class similarities, and background textures make this dataset challenging for distinguishing defects with similar appearances.

\textbf{GC10-DET:} The GC10-DET dataset \cite{Lv2020GC10}, collected from real industrial environments, comprises ten typical steel plate surface defect categories, covering diverse morphologies such as Punching Hole, Welding Line, Crescent Gap, Water Spot, Oil Spot, Silk Spot, Inclusion, Rolled Pit, Crease, and Waist Folding. It involves a richer variety of defect categories with larger variations in scale, orientation, and structural morphology, and is thus well-suited to verifying the adaptability of the model to industrial defects of multiple categories and morphologies.


\textbf{APDDD:} The APDDD dataset \cite{tianchi2018aluminium} is primarily dedicated to aluminum profile surface defect detection, consisting of 1,885 images that cover ten typical defect categories, including Depression, Non-conductivity, Scratch, Orange peel, Bottom leakage, Bruise, Pit, Convex powder, Coating crack, and Dirty spot. Owing to its distinct differences in substrate properties, surface textures, and defect imaging characteristics, it can be employed to construct detection scenarios involving material-domain discrepancies, thereby examining the performance stability of the model across different material scenarios.

\subsection{Evaluation Metrics}

Standard object detection metrics are adopted to quantify the detection performance. Precision is defined as
\begin{equation}
    \mathrm{Precision} = \frac{TP}{TP + FP}
\end{equation}
where $TP$ and $FP$ denote the number of true positives and false positives, respectively. Recall is expressed as
\begin{equation}
    \mathrm{Recall} = \frac{TP}{TP + FN}
\end{equation}
where $FN$ denotes the number of false negatives. The average precision (AP) is obtained as the area under the precision--recall curve for a single class,
\begin{equation}
    AP = \int_{0}^{1} P(R)\, dR
\end{equation}
and the mean average precision across all classes is given by
\begin{equation}
    mAP = \frac{1}{N} \sum_{i=1}^{N} AP_i
\end{equation}
where $N$ is the number of defect categories. In this study, results are primarily reported in terms of $mAP_{50}$ and $mAP_{50:95}$. $mAP_{50}$ corresponds to an IoU threshold of 0.5, whereas $mAP_{50:95}$ corresponds to the average $mAP$ over IoU thresholds ranging from 0.5 to 0.95 with a step size of 0.05. These metrics jointly capture classification accuracy, localization precision, and overall detection capability.


\subsection{Implementation Details}
\label{subsec:implementation}

All experiments were conducted on an NVIDIA GeForce RTX 3090 GPU
(24 GB) using Python 3.10 (https://www.python.org/),
PyTorch 2.5 (https://pytorch.org/), and CUDA 12.1
(https://developer.nvidia.com/cuda-toolkit).
All images were resized to $640 \times 640$, with a batch size
of 8 and a maximum of 300 training epochs.
FFRM-YOLO was trained using AdamW~\cite{Loshchilov2019AdamW}
with an initial learning rate of 0.002 and a weight decay
of 0.0005. Random flipping and rotation were applied for data
augmentation, and the auxiliary-box scaling ratio $\beta$
was set to 0.70.

Each dataset was divided into training and test sets
at a ratio of 7:3, with identical partitions used
across all methods. All architectural configurations,
hyperparameters, and model checkpoints were determined
using only the training data. The held-out test sets
were used exclusively for final performance evaluation
and were not involved in model development or selection. All comparison methods were re-trained using identical dataset partitions, input resolution, optimizer, learning rate schedule, and evaluation metrics to maintain a consistent experimental protocol.
To assess result variability, YOLO11n, WaveMamba-YOLO, and
FFRM-YOLO were each trained using three random seeds
(0, 1, and 42). Precision and Recall are reported as
the means of the three runs, while mAP$_{50}$ and
mAP$_{50:95}$ are presented as mean $\pm$ standard
deviation. The same repeated-run protocol was applied
to the baseline, GLRM--MDRM, and complete configurations
in the component ablation study; other ablations
report individual-run results. Results for the remaining
comparison methods are based on single runs.

Model complexity was evaluated using parameter counts and
FLOPs at an input resolution of $640 \times 640$.
Inference throughput (FPS) was measured on the same RTX 3090
GPU with batch size 1 using PyTorch FP32 inference, including
preprocessing, model inference, and post-processing.


% =========================================================
% 4.4 Comparative Experiments
% =========================================================
\subsection{Comparative Experiments}
\label{subsec:comparison}

We compare FFRM-YOLO with representative YOLO-based detectors and recent defect detection methods on NEU-DET, GC10-DET, and APDDD under the evaluation protocol described in Section~\ref{subsec:implementation}. The following subsections present the results for each dataset, followed by a joint analysis of detection accuracy and model complexity.

\begin{table}[!t]
\caption{Detection performance and model complexity on NEU-DET. \textbf{Bold} and \underline{underlined} values indicate the best and second-best detection results, respectively. Model complexity and FPS are evaluated at an input resolution of $640\times640$. FPS is measured on a single NVIDIA RTX 3090 GPU with batch size~1 using PyTorch FP32 end-to-end inference(including pre- and post-processing).}
\label{tab:comparison_neu}
\centering
\footnotesize
\setlength{\tabcolsep}{1.3pt}
\renewcommand{\arraystretch}{1.12}
\begin{tabular*}{\textwidth}{@{\extracolsep{\fill}}llccccccc@{}}
\toprule
Method & Published & P & R & $\mathrm{mAP}_{50}$ & $\mathrm{mAP}_{50:95}$ & Params (M) & FLOPs (G) & FPS \\
\midrule
YOLOv5n~\cite{Jocher2020YOLOv5} & Ultralytics'20 & 73.2 & 71.3 & 74.8 & 41.4 & 2.5 & 7.1 &\underline{168.6}\\
YOLOv8n~\cite{Jocher2023YOLOv8} & Ultralytics'23 & 74.4 & 73.9 & 75.6 & 42.2 & 3.0 & 8.1 &145.2\\
YOLOv9t~\cite{Wang2024YOLOv9} & ECCV'24 & 75.5 & 68.0 & 76.6 & 43.3 & \textbf{2.0} & 7.6 &155.4\\
YOLOv10n~\cite{Wang2024YOLOv10} & NeurIPS'24 & 72.1 & 67.2 & 73.1 & 40.5 & \underline{2.3} & \underline{6.5} &\textbf{175.5}\\
YOLO11n~\cite{Jocher2024YOLO11} & Ultralytics'24 & 73.5 & 73.6 & 75.7\(\pm\)0.8 & 43.0\(\pm\)0.3 & 2.6 & \textbf{6.3} &165.8\\
ES-Net~\cite{Yu2022ESNet} & TIM'22 & 74.8 & 69.9 & 76.4 & 42.0 & 2.9 & 8.2 &111.8\\
GC-Net~\cite{Liu2025GCNet} & PR'25 & 73.1 & 71.1 & 76.7 & 42.7 & 3.1 & 8.7 &100.9\\
MD-YOLO~\cite{Zheng2024MDYOLO} & OLE'24 & 72.7 & 72.4 & 74.9 & 41.2 & 2.5 & 13.4 &85.6\\
SDCM-YOLO~\cite{Deng2026SDCMYOLO} & ERE'26 & \underline{76.7} & 67.2 & 76.4 & 41.8 & 4.2 & 12.4 &90.3 \\
MWYOLO~\cite{Chen2026MWYOLO} & AIE'26 & 74.5 & 71.8 & 77.3 & 42.9 & 3.8 & 11.2 &103.5\\
WaveMamba-YOLO~\cite{Ge2026WaveMambaYOLO} & PLOS ONE'26 & 74.2 &\underline{74.3} & 77.1\(\pm\)0.6 & 43.3\(\pm\)0.4 & 3.3 & 6.9 & 95.5 \\
SH-DETR~\cite{Wu2025SHDETR} & PONE'25 & \textbf{90.7} & \textbf{78.1} & \textbf{81.3} & \underline{43.6} & 16.4 & 44.3 &65.4\\
\midrule
\textbf{FFRM-YOLO (Ours)} & -- & 74.4 & 74.1 & \underline{78.6\(\pm\)0.5} & \textbf{44.5\(\pm\)0.6} & 4.8 & 7.7 &135.2\\
\bottomrule
\end{tabular*}
\end{table}



\begin{table}[!t]
\caption{Detection performance and model complexity on GC10-DET. \textbf{Bold} and \underline{underlined} values indicate the best and second-best detection results, respectively. Model complexity and FPS are evaluated at an input resolution of $640\times640$. FPS is measured on a single NVIDIA RTX 3090 GPU with batch size~1 using PyTorch FP32 end-to-end inference(including pre- and post-processing).}
\label{tab:comparison_gc10}
\centering
\footnotesize
\setlength{\tabcolsep}{1.3pt}
\renewcommand{\arraystretch}{1.12}
\begin{tabular*}{\textwidth}{@{\extracolsep{\fill}}llccccccc@{}}
\toprule
Method & Published & P & R & $\mathrm{mAP}_{50}$ & $\mathrm{mAP}_{50:95}$ & Params (M) & FLOPs (G) & FPS \\
\midrule
YOLOv5n~\cite{Jocher2020YOLOv5} & Ultralytics'20 & 65.2 & 59.3 & 61.2 & 30.0 & 2.5 & 7.1 & \underline{168.6} \\
YOLOv8n~\cite{Jocher2023YOLOv8} & Ultralytics'23 & 59.0 & 63.5 & 61.8 & \underline{32.6} & 3.0 & 8.1 & 145.2\\
YOLOv9t~\cite{Wang2024YOLOv9} & ECCV'24 & 65.9 & 59.3 & 61.1 & 29.9 & \textbf{2.0} & 7.6 & 155.4 \\
YOLOv10n~\cite{Wang2024YOLOv10} & NeurIPS'24 & 57.5 & 52.7 & 54.6 & 28.5 & \underline{2.3} & \underline{6.5} & \textbf{175.5}\\
YOLO11n~\cite{Jocher2024YOLO11} & Ultralytics'24 & 68.3 & 63.1 & 63.4\(\pm\)0.6 & 31.2\(\pm\)0.5 & 2.6 & \textbf{6.3} & 165.8\\
ES-Net~\cite{Yu2022ESNet} & TIM'22 & 60.4 & 58.3 & 59.2 & 29.6 & 2.9 & 8.2 & 111.8\\
GC-Net~\cite{Liu2025GCNet} & PR'25 & 64.4 & 60.7 & 62.2 & 30.6 & 3.1 & 8.7 & 100.9\\
MD-YOLO~\cite{Zheng2024MDYOLO} & OLE'24 & 55.4 & 56.6 & 54.8 & 29.4 & 2.5 & 13.4 &85.6\\
SDCM-YOLO~\cite{Deng2026SDCMYOLO} & ERE'26 & 65.2 & 58.9 & 61.9 & 29.9 & 4.2 & 12.4 & 90.3\\
MWYOLO~\cite{Chen2026MWYOLO} & AIE'26 & 66.7 & 59.4 & 60.8 & 30.0 & 3.8 & 11.2 & 103.5\\
WaveMamba-YOLO~\cite{Ge2026WaveMambaYOLO} & PLOS ONE'26 & 68.7 & 57.6 &63.9\(\pm\)0.2 & 31.9\(\pm\)0.3 & 3.3 & 6.9 & 95.5\\
SH-DETR~\cite{Wu2025SHDETR} & PONE'25 & \textbf{75.4} & \textbf{61.2} & 63.8 & 32.2 & 16.4 & 44.3 &65.4\\
\midrule
    \textbf{FFRM-YOLO (Ours)} & -- & \underline{71.2} & \textbf{63.7} & \textbf{66.1\(\pm\)0.3} & \textbf{34.5\(\pm\)0.4} & 4.8 & 7.7 & 135.2\\
\bottomrule
\end{tabular*}
\end{table}



\noindent\textbf{Results on NEU-DET:} 
Table~\ref{tab:comparison_neu} presents the detection results on NEU-DET. FFRM-YOLO achieves an $\mathrm{mAP}_{50}$ of $(78.6 \pm 0.5)\%$ and an $\mathrm{mAP}_{50:95}$ of $(44.5 \pm 0.6)\%$, outperforming YOLO11n by 2.9 and 1.5 percentage points in the corresponding mean values, respectively. Precision and recall also increase from 73.5\% and 73.6\% to 74.4\% and 74.1\%. Compared with WaveMamba-YOLO, a closely related frequency-domain method, FFRM-YOLO improves the mean $\mathrm{mAP}_{50}$ and $\mathrm{mAP}_{50:95}$ by 1.5 and 1.2 percentage points, respectively. Although SH-DETR achieves a higher $\mathrm{mAP}_{50}$ (81.3\% vs. 78.6\%), FFRM-YOLO obtains the highest mean $\mathrm{mAP}_{50:95}$ among the evaluated methods (44.5\%), while requiring substantially fewer parameters and FLOPs. These results indicate competitive detection accuracy, particularly under stricter localization thresholds.



\noindent\textbf{Results on GC10-DET:} The detection results on GC10-DET are reported in Table~\ref{tab:comparison_gc10}. FFRM-YOLO achieves an $\mathrm{mAP}_{50}$ of $(66.1 \pm 0.3)\%$ and an $\mathrm{mAP}_{50:95}$ of $(34.5 \pm 0.4)\%$, outperforming YOLO11n by 2.7 and 3.3 percentage points in the corresponding mean values, respectively. Precision and recall also increase from 68.3\% and 63.1\% to 71.2\% and 63.7\%, respectively. Compared with WaveMamba-YOLO, FFRM-YOLO improves the mean $\mathrm{mAP}_{50}$ and $\mathrm{mAP}_{50:95}$ by 2.2 and 2.6 percentage points, respectively. It also exceeds SH-DETR by 2.3 percentage points in both mAP metrics, although SH-DETR achieves higher precision. Overall, FFRM-YOLO achieves the highest mean values for both mAP metrics among the evaluated methods, indicating its effectiveness in detecting defects with diverse categories and appearances.



\begin{table}[!t]
\caption{Detection performance and model complexity on APDDD. \textbf{Bold} and \underline{underlined} values indicate the best and second-best detection results, respectively. Dashes indicate unavailable results. Model complexity and FPS are evaluated at an input resolution of $640\times640$. FPS is measured on a single NVIDIA RTX 3090 GPU with batch size~1 using PyTorch FP32 end-to-end inference(including pre- and post-processing).}
\label{tab:comparison_apddd}
\centering
\footnotesize
\setlength{\tabcolsep}{1.3pt}
\renewcommand{\arraystretch}{1.12}
\begin{tabular*}{\textwidth}{@{\extracolsep{\fill}}llccccccc@{}}
\toprule
Method & Published & P & R & $\mathrm{mAP}_{50}$ & $\mathrm{mAP}_{50:95}$ & Params (M) & FLOPs (G) & FPS \\
\midrule
YOLOv5n~\cite{Jocher2020YOLOv5} & Ultralytics'20 & 72.9 & 50.5 & 56.5 & 32.6 & 2.5 & 7.1 &\underline{168.6}\\
YOLOv8n~\cite{Jocher2023YOLOv8} & Ultralytics'23 & \underline{70.0} & 53.5 & 57.3 & 33.7 & 3.0 & 8.1 &145.2\\
YOLOv9t~\cite{Wang2024YOLOv9} & ECCV'24 & 66.3 & 51.1 & 53.6 & 31.9 & \textbf{2.0} & 7.6 &155.4\\
YOLOv10n~\cite{Wang2024YOLOv10} & NeurIPS'24 & 54.0 & 46.6 & 46.0 & 27.3 & \underline{2.3} & \underline{6.5} &\textbf{175.5}\\
YOLO11n~\cite{Jocher2024YOLO11} & Ultralytics'24 & 68.6 & 55.5 & 57.6\(\pm\)0.2 & \underline{34.2\(\pm\) 0.3} & 2.6 & \textbf{6.3} &165.8 \\
ES-Net~\cite{Yu2022ESNet} & TIM'22 & 62.5 & 54.4 & 54.8 & 31.7 & 2.9 & 8.2 &111.8\\
GC-Net~\cite{Liu2025GCNet} & PR'25 & 68.5 & 55.1 & 57.5 & 33.4 & 3.1 & 8.7 &100.9\\
MD-YOLO~\cite{Zheng2024MDYOLO} & OLE'24 & 75.7 & 52.3 & 56.3 & 33.2 & 2.5 & 13.4 &85.6\\
SDCM-YOLO~\cite{Deng2026SDCMYOLO} & ERE'26 & 65.8 & 53.6 & 56.1 & 31.4 & 4.2 & 12.4 &90.3\\
MWYOLO~\cite{Chen2026MWYOLO} & AIE'26 & 65.5 & \underline{56.5} & 55.8 & 32.3 & 3.8 & 11.2 &103.5 \\
WaveMamba-YOLO~\cite{Ge2026WaveMambaYOLO} & PLOS ONE'26 & 67.4 & 52.1 & 57.6\(\pm\)0.3 & 32.4\(\pm\)0.2 & 3.3 & 6.9 & 95.5\\
SH-DETR~\cite{Wu2025SHDETR} & PONE'25 & 69.7 & 55.4 & \underline{58.6} & 33.9 & 16.4 & 44.3 &65.4\\
\midrule
\textbf{FFRM-YOLO (Ours)} & -- & \textbf{74.1} & \textbf{56.8} & \textbf{60.6\(\pm\)0.6} & \textbf{35.3\(\pm\)0.4} & 4.8 & 7.7 & 135.2\\
\bottomrule
\end{tabular*}
\end{table}



\noindent\textbf{Results on APDDD:} Table~\ref{tab:comparison_apddd} presents the detection results on APDDD. FFRM-YOLO achieves an $\mathrm{mAP}_{50}$ of $(60.6 \pm 0.6)\%$ and an $\mathrm{mAP}_{50:95}$ of $(35.3 \pm 0.4)\%$, outperforming YOLO11n by 3.0 and 1.1 percentage points in the corresponding mean values, respectively. Precision and recall also increase from 68.6\% and 55.5\% to 74.1\% and 56.8\%, respectively. Compared with WaveMamba-YOLO, FFRM-YOLO improves the mean $\mathrm{mAP}_{50}$ and $\mathrm{mAP}_{50:95}$ by 3.0 and 2.9 percentage points, respectively. It also outperforms SH-DETR by 2.0 and 1.4 percentage points in the two mAP metrics. Among the evaluated methods, FFRM-YOLO achieves the highest recall and both mAP metrics, although MD-YOLO obtains a higher precision of 75.7\%. These results further demonstrate the effectiveness of the proposed model across different defect datasets.



\noindent\textbf{Accuracy and Model Complexity:}Tables~\ref{tab:comparison_neu}--\ref{tab:comparison_apddd} summarize the detection performance and model complexity at an input resolution of $640 \times 640$. FFRM-YOLO requires 4.8~M parameters and 7.7~GFLOPs, representing increases of 2.2~M parameters and 1.4~GFLOPs over YOLO11n, respectively. These additional costs are accompanied by mean $\mathrm{mAP}_{50}$ improvements of 2.9, 2.7, and 3.0 percentage points on NEU-DET, GC10-DET, and APDDD, respectively. Compared with SH-DETR, FFRM-YOLO uses 70.7\% fewer parameters and 82.6\% fewer FLOPs, while achieving higher $\mathrm{mAP}_{50:95}$ on all three datasets. It also reduces FLOPs by 31.3\% relative to MWYOLO, despite requiring 1.0~M additional parameters. Under the same inference setting, FFRM-YOLO achieves 135.2~FPS, compared with 165.8~FPS for YOLO11n and 95.5~FPS for WaveMamba-YOLO. Overall, these results indicate a competitive accuracy--efficiency trade-off, although the proposed model incurs additional computational overhead relative to YOLO11n.

\FloatBarrier

\subsection{Ablation Studies}

\noindent\textbf{Ablation Study of Components.} We evaluate eight configurations of GLRM, MDRM, and Inner-CIoU on NEU-DET and GC10-DET, all based on YOLO11n with consistent training settings except for the components under investigation. The mAP metrics for variants~1, 5, and~8 are reported as mean $\pm$ standard deviation over three independent runs.

\begin{table}[!t]
    \centering
    \caption{Component ablation results on NEU-DET and GC10-DET.
    All variants are based on YOLO11n.
    \textbf{Bold} values indicate the best result for each metric within each dataset.}
    \label{tab:component_ablation}

    \fontsize{6.0pt}{7.5pt}\selectfont
    \setlength{\tabcolsep}{0.4pt}
    \renewcommand{\arraystretch}{1.08}

    \begin{tabular*}{\textwidth}
    {@{\extracolsep{\fill}}cccc*{8}{c}@{}}
    \toprule

    \multirow{2}{*}{\textbf{No.}} &
    \multirow{2}{*}{\textbf{GLRM}} &
    \multirow{2}{*}{\textbf{MDRM}} &
    \multirow{2}{*}{\textbf{Inner-CIoU}} &
    \multicolumn{4}{c}{\textbf{NEU-DET}} &
    \multicolumn{4}{c}{\textbf{GC10-DET}} \\

    \cmidrule(lr){5-8}
    \cmidrule(lr){9-12}

    & & & &
    \textbf{P} &
    \textbf{R} &
    \textbf{mAP$_{50}$} &
    \textbf{mAP$_{50:95}$} &
    \textbf{P} &
    \textbf{R} &
    \textbf{mAP$_{50}$} &
    \textbf{mAP$_{50:95}$} \\

    \midrule

    1 & -- & -- & -- &
    73.5 & 73.6 & 75.7\(\pm\)0.8 & 43.0\(\pm\)0.3 &
    68.3 & 63.1 & 63.4\(\pm\)0.6 & 31.2\(\pm\)0.5\\

    2 & $\checkmark$ & -- & -- &
    73.4 & \textbf{74.9} & 76.3 & 43.2 &
    68.7 & 63.5 & 64.7 & 32.4 \\

    3 & -- & $\checkmark$ & -- &
    73.6 & 74.7 & 76.6 & 43.8 &
    69.6 & \textbf{64.8} & 65.2 & 32.6 \\

    4 & -- & -- & $\checkmark$ &
    72.4 & 73.7 & 75.3 & 42.6 &
    69.0 & 62.8 & 64.5 & 31.5 \\

    5 & $\checkmark$ & $\checkmark$ & -- &
    74.3 & 73.9 & 78.2\(\pm\)0.3 & 44.0\(\pm\)0.2 &
    70.9 & 63.4 & 65.7\(\pm\)0.3 & 34.1\(\pm\)0.3 \\

    6 & $\checkmark$ & -- & $\checkmark$ &
    73.1 & 73.8 & 77.2 & 43.4 &
    69.5 & 63.5 & 65.1 & 32.1 \\

    7 & -- & $\checkmark$ & $\checkmark$ &
    73.8 & 73.6 & 77.9 & 44.2 &
    70.5 & 64.4 & 65.6 & 33.9 \\

    \midrule

    8 & $\checkmark$ & $\checkmark$ & $\checkmark$ &
    \textbf{74.4} & 74.1 &
    \textbf{78.6\(\pm\)0.5} & \textbf{44.5\(\pm\)0.6} &
    \textbf{71.2} & 63.7 &
    \textbf{66.1\(\pm\)0.3} & \textbf{34.5\(\pm\)0.4} \\

    \botrule
    \end{tabular*}
\end{table}


As shown in Table~\ref{tab:component_ablation}, variants~2 and~3 improve both mAP metrics over the baseline (variant~1) on NEU-DET and GC10-DET, supporting the individual contributions of GLRM and MDRM. Their combination (variant~5) further improves performance, achieving $\mathrm{mAP}_{50}$ values of $(78.2 \pm 0.3)\%$ and $(65.7 \pm 0.3)\%$, and $\mathrm{mAP}_{50:95}$ values of $(44.0 \pm 0.2)\%$ and $(34.1 \pm 0.3)\%$ on the two datasets, respectively. These correspond to mean $\mathrm{mAP}_{50}$ improvements of 2.5 and 2.3 percentage points over the baseline, demonstrating the complementary contributions of both modules.

Inner-CIoU alone (variant~4) yields mixed effects across the two datasets, while its combination with individual modules (variants~6 and~7) produces configuration-dependent improvements. Adding Inner-CIoU to GLRM and MDRM (variant~8) increases mean $\mathrm{mAP}_{50}$ by 0.4 percentage points on both datasets and mean $\mathrm{mAP}_{50:95}$ by 0.5 and 0.4 percentage points, respectively. The complete model achieves $\mathrm{mAP}_{50}$ values of $(78.6 \pm 0.5)\%$ and $(66.1 \pm 0.3)\%$, and $\mathrm{mAP}_{50:95}$ values of $(44.5 \pm 0.6)\%$ and $(34.5 \pm 0.4)\%$, representing the best mAP performance among the evaluated variants. These results indicate that GLRM and MDRM provide the primary gains, with Inner-CIoU offering modest additional improvements.



\noindent\textbf{Ablation Study of FFRM Placement.}
We compare six placement configurations to examine how the distribution of FFRM affects detection performance: no replacement, two partial replacement configurations, backbone-only replacement, neck-only replacement, and replacement throughout both the backbone and neck. The replacement sites are identified using the block labels in Fig.~\ref{fig:architecture}, where B1--B4 denote backbone blocks and N1--N4 denote neck blocks. All other experimental settings are kept unchanged.

\begin{table}[!t]
    \caption{Effect of FFRM placement. B1--B4 and N1--N4 denote the backbone and neck blocks labeled in Fig.~\ref{fig:architecture}, respectively. Inner-CIoU is used for all configurations. Paired values correspond to NEU-DET/GC10-DET. \textbf{Bold} values indicate the best results for each dataset.}
    \label{tab:placement}
    \centering
    \footnotesize
    \setlength{\tabcolsep}{3pt}
    \renewcommand{\arraystretch}{1.1}
    \begin{tabular*}{\textwidth}
    {@{\extracolsep{\fill}}lcccc@{}}
    \toprule
    Replacement Sites &
    Precision &
    Recall &
    $\mathrm{mAP}_{50}$ &
    $\mathrm{mAP}_{50:95}$ \\
    \midrule
    None
    & 72.4/69.0
    & 73.7/62.8
    & 75.3/64.5
    & 42.6/31.5 \\

    B3, B4, N1, N2
    & \textbf{75.2}/69.0
    & 66.4/59.6
    & 75.3/63.8
    & 41.9/32.3 \\

    B4, N1--N4
    & 69.4/64.1
    & 72.0/57.3
    & 75.6/61.6
    & 41.9/32.0 \\

    Backbone (B1--B4)
    & 72.3/69.7
    & 72.1/63.4
    & 77.6/64.5
    & 43.4/33.7 \\

    Neck (N1--N4)
    & 72.8/69.4
    & \textbf{74.5}/63.3
    & 77.9/64.1
    & 43.6/33.2 \\

    All (Ours)
    & 74.4/\textbf{71.2}
    & 74.1/\textbf{63.7}
    & \textbf{78.8}/\textbf{66.3}
    & \textbf{44.7}/\textbf{34.6} \\
    \bottomrule
    \end{tabular*}
\end{table}

\begin{table}[!t]
\caption{Ablation study of model capacity and frequency decomposition
on NEU-DET and GC10-DET. Paired values correspond to
NEU-DET/GC10-DET. The capacity-matched YOLO11n increases
model capacity by adjusting the C3k2 configurations, while
the No DWT/IDWT variant replaces wavelet-based processing
with conventional spatial-domain operations. Bold values
indicate the best detection results. Parameters and FLOPs
are measured at an input resolution of $640\times640$.}
\label{tab:ffrm_dwt}
\centering
\footnotesize
\setlength{\tabcolsep}{1.8pt}
\renewcommand{\arraystretch}{1.1}
\begin{tabular*}{\columnwidth}{@{\extracolsep{\fill}}lcccccccc@{}}
\toprule
Variant & Precision & Recall & $\mathrm{mAP}_{50}$ & $\mathrm{mAP}_{50:95}$&Params (M)& FLOPs (G) \\
\midrule
Baseline  &73.0/68.3& 73.6/63.1 & 75.9/63.9 & 42.9/31.8 & \textbf{2.6} & \textbf{6.3} \\
YOLO11n(Capacity-matched)   &74.1/66.2  &69.7/62.6  &76.5/62.8 &43.5/30.9  & 4.5 & 7.8\\
FFRM-YOLO(No DWT/IDWT)  &73.8/69.4  &71.3/63.0   &77.5/63.7  & 43.2/32.5 & 4.9 & 7.8\\
FFRM-YOLO(Full)    & \textbf{74.4/71.2} & \textbf{74.1/63.7} & \textbf{78.8/66.3} & \textbf{44.7/34.6} & 4.8 & 7.7\\
\bottomrule
\end{tabular*}
\end{table}


As shown in Table~\ref{tab:placement}, the two partial replacement configurations produce mixed changes relative to the ``None'' configuration. On NEU-DET, they maintain or slightly improve $\mathrm{mAP}_{50}$ but reduce $\mathrm{mAP}_{50:95}$, whereas on GC10-DET, they decrease $\mathrm{mAP}_{50}$ while slightly improving $\mathrm{mAP}_{50:95}$. Replacing all eligible blocks within either the backbone or the neck generally yields stronger performance, particularly on NEU-DET, while the gains on GC10-DET remain metric-dependent. In the single-run placement comparison, replacing blocks in both the backbone and neck achieves the highest mAP values among the evaluated configurations, reaching 78.8\%/44.7\% on NEU-DET and 66.3\%/34.6\% on GC10-DET for mAP\(_{50}\)/mAP\(_{50:95}\), respectively. Although this configuration does not achieve the highest precision or recall in every case, its higher mAP values suggest that applying frequency refinement during both feature extraction and multi-scale fusion is beneficial under the evaluated setting. We therefore adopt replacement of all eligible C3k2 blocks in the backbone and neck as the default configuration.

\noindent\textbf{Ablation on Model Capacity and Frequency Decomposition.} To examine whether the performance gains arise from increased model capacity, we construct a capacity-matched YOLO11n variant by setting the repetition parameter to 6 and the expansion ratio to 0.625 for the C3k2 modules at backbone layer 8 and head layer 22 in the original YOLO11n YAML configuration, while keeping all other settings unchanged. As shown in Table~\ref{tab:ffrm_dwt}, this variant has 4.5~M parameters and 7.8~GFLOPs, comparable to FFRM-YOLO (4.8~M and 7.7~GFLOPs). However, FFRM-YOLO achieves higher $\mathrm{mAP}{50}$ by 2.3 and 3.5 percentage points and higher $\mathrm{mAP}{50:95}$ by 1.2 and 3.7 percentage points on NEU-DET and GC10-DET, respectively, suggesting that increased model capacity alone does not explain the performance gains.

To assess the contribution of frequency decomposition, we replace DWT/IDWT with conventional spatial-domain processing while retaining the GLRM--MDRM refinement structure. Compared with this variant, FFRM-YOLO improves $\mathrm{mAP}{50}$ by 1.3 and 2.6 percentage points and $\mathrm{mAP}{50:95}$ by 1.5 and 2.1 percentage points on NEU-DET and GC10-DET, respectively, despite having slightly fewer parameters and FLOPs (4.8~M/7.7~G versus 4.9~M/7.8~G). These results support the contribution of wavelet-based decomposition and reconstruction beyond conventional spatial-domain processing.




\noindent\textbf{Ablation Study of GLRM.} We evaluate the local branch in GLRM by comparing its removal with depthwise convolutions using $3\times3$, $5\times5$, and $7\times7$ kernels, while keeping the remaining configuration unchanged. As shown in Table~\ref{tab:glrm_kernel}, the $7\times7$ configuration achieves the highest values for both mAP metrics on both datasets. Compared with removing the local branch, it improves $\mathrm{mAP}_{50}$ and $\mathrm{mAP}_{50:95}$ by 2.5\% and 2.9\% on NEU-DET, and by 2.5\% and 1.0\% on GC10-DET. However, increasing the kernel size does not improve every metric monotonically, and the $7\times7$ configuration does not achieve the highest precision or recall in every case. Relative to the $5\times5$ configuration, it improves the two mAP metrics by 1.3\% and 1.0\% on NEU-DET and by 0.8\% and 1.3\% on GC10-DET, with an additional 0.2 GFLOPs. These results support using local convolution to complement SS2D-based global modeling and selecting $7\times7$ as the default kernel among the tested configurations.

\begin{table}[!t]
\caption{Effect of local-branch kernel sizes in GLRM. Paired values correspond to NEU-DET/GC10-DET. $0\times0$ denotes removal of the local branch. \textbf{Bold} values indicate the best results.}
    \label{tab:glrm_kernel}
    \centering
    \footnotesize
    \setlength{\tabcolsep}{3pt}
    \renewcommand{\arraystretch}{1.1}
    \begin{tabular*}{\textwidth}
    {@{\extracolsep{\fill}}lcccccc@{}}
    \toprule
    Size &
    Precision &
    Recall &
    $\mathrm{mAP}_{50}$ &
    $\mathrm{mAP}_{50:95}$ &
    Params &
    GFLOPs\\
    \midrule
    $0\times0$
    & 73.1/\textbf{72.6}
    & 74.3/58.3
    & 76.3/63.8
    & 41.8/33.6
    & \textbf{4.5}
    & \textbf{7.1} \\

    $3\times3$
    & 73.5/70.6
    & 74.2/59.7
    & 77.0/64.3
    & 42.3/32.5
    & 4.7
    & 7.4 \\

    $5\times5$
    & 74.1/70.3
    & \textbf{74.6}/62.0
    & 77.5/65.5
    & 43.7/33.3
    & 4.7
    & 7.5 \\

    $7\times7$
    & \textbf{74.4}/71.2
    & 74.1/\textbf{63.7}
    & \textbf{78.8}/\textbf{66.3}
    & \textbf{44.7}/\textbf{34.6}
    & 4.8
    & 7.7 \\
    \bottomrule
    \end{tabular*}
\end{table}

\begin{table}[!t]
    \centering
    \caption{Effect of kernel allocation in MK-DeConv. LH, HL, and HH specify the convolution kernel sizes for the corresponding subbands. Paired results correspond to NEU-DET/GC10-DET. \textbf{Bold} values indicate the best results.}
    \label{tab:mdrm_kernel}
    \footnotesize
    \setlength{\tabcolsep}{3pt}
    \renewcommand{\arraystretch}{1.15}
    \begin{tabular*}{\textwidth}
    {@{\extracolsep{\fill}}cccccccc@{}}
    \toprule
    \textbf{No.} &
    \textbf{LH} &
    \textbf{HL} &
    \textbf{HH} &
    \textbf{Precision} &
    \textbf{Recall} &
    \textbf{mAP$_{50}$} &
    \textbf{mAP$_{50:95}$} \\
    \midrule
    1 & $3\times3$ & $3\times3$ & $3\times3$
    & 70.2/64.5
    & 72.2/58.0
    & 76.1/61.9
    & 42.0/32.4 \\

    2 & $5\times5$ & $5\times5$ & $5\times5$
    & \textbf{74.4}/65.6
    & 69.3/57.6
    & 75.6/62.1
    & 41.7/31.9 \\

    3 & $7\times7$ & $7\times7$ & $7\times7$
    & 73.8/67.5
    & 70.0/60.3
    & 75.9/63.5
    & 41.9/32.1 \\

    4 & $7\times7$ & $5\times5$ & $3\times3$
    & 68.1/70.4
    & 73.2/60.8
    & 75.7/65.1
    & 42.0/33.3 \\
    \midrule
    5 & $3\times3$ & $5\times5$ & $7\times7$
    & \textbf{74.4}/\textbf{71.2}
    & \textbf{74.1}/\textbf{63.7}
    & \textbf{78.8}/\textbf{66.3}
    & \textbf{44.7}/\textbf{34.6} \\
    \bottomrule
    \end{tabular*}
\end{table}

\noindent\textbf{Ablation Study of MDRM.} 
We compare five kernel allocation strategies in MK-DeConv while
keeping the remaining configuration unchanged. Configurations~1--3
use uniform kernel sizes, whereas configurations~4 and~5 assign
different kernels to the LH, HL, and HH subbands. As shown in
Table~\ref{tab:mdrm_kernel}, configuration~5, which assigns
$3\times3$, $5\times5$, and $7\times7$ kernels to LH, HL, and HH,
respectively, achieves the highest mAP values on both datasets.
Compared with the reversed allocation (configuration~4), it
improves $\mathrm{mAP}_{50}$ and $\mathrm{mAP}_{50:95}$ by
3.1 and 2.7 percentage points on NEU-DET, and by 1.2 and
1.3 percentage points on GC10-DET, respectively. Since both
configurations use the same kernel sizes, these differences
suggest that subband-specific receptive-field allocation is
important. A possible explanation is that smaller kernels
better preserve localized directional details in LH, while
larger kernels provide broader contextual integration for
the more complex high-frequency patterns in HH, with HL
using an intermediate receptive field. Although this
interpretation requires further investigation, the results
support selecting configuration~5 among the evaluated
strategies.



\begin{table}[!t]
    \centering
    \caption{Comparison of bounding-box regression losses. Paired values correspond to NEU-DET/GC10-DET. \textbf{Bold} values indicate the best results.}
    \label{tab:loss}
    \footnotesize
    \setlength{\tabcolsep}{3pt}
    \renewcommand{\arraystretch}{1.15}
    \begin{tabular*}{\textwidth}
    {@{\extracolsep{\fill}}lcccc@{}}
    \toprule
    \textbf{Loss} &
    \textbf{Precision} &
    \textbf{Recall} &
    \textbf{mAP$_{50}$} &
    \textbf{mAP$_{50:95}$} \\
    \midrule
    CIoU~\cite{Zheng2020distance}
    & 74.3/70.9
    & 73.9/63.4
    & 78.1/65.8
    & 44.5/34.4 \\

    Inner-EIoU~\cite{Zhang2023InnerIoU}
    & 73.4/69.2
    & 72.7/58.0
    & 77.2/62.6
    & 42.5/31.7 \\

    Inner-DIoU~\cite{Zhang2023InnerIoU}
    & 72.9/68.6
    & 72.1/60.4
    & 77.0/63.7
    & 42.2/32.5 \\

    Inner-GIoU~\cite{Zhang2023InnerIoU}
    & 71.6/66.1
    & 71.0/59.3
    & 75.7/62.6
    & 41.9/31.7 \\

    Inner-SIoU~\cite{Zhang2023InnerIoU}
    & 72.9/66.5
    & 70.1/61.8
    & 76.5/64.0
    & 42.2/32.6 \\
    \midrule
    Inner-CIoU~\cite{Zhang2023InnerIoU}
    & \textbf{74.4}/\textbf{71.2}
    & \textbf{74.1}/\textbf{63.7}
    & \textbf{78.8}/\textbf{66.3}
    & \textbf{44.7}/\textbf{34.6} \\
    \bottomrule
    \end{tabular*}
\end{table}

\noindent\textbf{Ablation Study of Loss Functions.}
We compare CIoU with several Inner-IoU variants to evaluate the effect of the regression loss on defect detection, while keeping the feature extraction architecture and other training settings unchanged. As shown in Table~\ref{tab:loss}, Inner-CIoU achieves the highest precision, recall, and both mAP metrics on NEU-DET and GC10-DET. Compared with CIoU, it improves $\mathrm{mAP}_{50}$ and $\mathrm{mAP}_{50:95}$ by 0.7\% and 0.2\% on NEU-DET, and by 0.5\% and 0.2\% on GC10-DET, respectively. These results show modest additional gains when auxiliary-box overlap is combined with the center-distance and aspect-ratio penalties of CIoU. In contrast, all other tested Inner-IoU variants yield lower values for both mAP metrics than CIoU on both datasets, indicating that auxiliary-box overlap alone does not guarantee improved performance. Among the evaluated losses, Inner-CIoU provides the best results under the present experimental settings and is therefore selected as the regression loss of FFRM-YOLO.


Overall, the ablation studies support the effectiveness of the main components and design choices of FFRM-YOLO. GLRM and MDRM provide complementary refinement of low-frequency structural information and high-frequency texture details, respectively, while Inner-CIoU offers modest additional improvements when combined with both modules. Furthermore, the capacity-matched and spatial-domain control experiments support the contribution of wavelet-based frequency decomposition beyond increased model capacity. These findings highlight the complementary roles of frequency decomposition, feature refinement, and localization optimization.



\FloatBarrier


\subsection{Frequency-Feature Visualization}


\begin{figure}[!t]
    \centering
    \includegraphics[width=0.98\textwidth]{idwt.png}
\caption{Frequency-feature visualization for five defect categories on NEU-DET. The first column shows input images with ground-truth bounding boxes (GT). The remaining columns show the refined low-frequency output (LL) and the three high-frequency output groups (LH, HL, and HH) before IDWT reconstruction. The high-frequency labels indicate their respective reconstruction positions after channel interaction in MDRM.}
    \label{fig:subband_vis}
\end{figure}

Fig.~\ref{fig:subband_vis} visualizes the refined subband features $Y_{LL}$, $Y_{LH}$, $Y_{HL}$, and $Y_{HH}$ before IDWT reconstruction in the shallowest FFRM (B1 in Fig.~\ref{fig:architecture}, stride 4), where fine-grained spatial details are preserved. The low-frequency output $Y_{LL}$ is generated by GLRM, while the three high-frequency outputs are produced by MDRM. 




\begin{table}[!t]
\centering
\caption{Per-class detection performance of YOLO11n and
FFRM-YOLO on NEU-DET. Bold values indicate the
best result for each defect category.}
\label{tab:per_class_neu}

\small
\setlength{\tabcolsep}{4pt}
\renewcommand{\arraystretch}{1.08}

\begin{tabular*}{\textwidth}{@{\extracolsep{\fill}}lcccccc@{}}
\hline
\noalign{\vskip 2pt}

\raisebox{-1.1ex}{Method}
& Crazing
& Inclusion
& Patches
& \shortstack{Pitted\\Surface}
& \shortstack{Rolled-in\\Scale}
& Scratches \\

\noalign{\vskip -8pt}
& \multicolumn{6}{@{}c@{}}{\hrulefill} \\
\noalign{\vskip 0pt}

& mAP$_{50}$
& mAP$_{50}$
& mAP$_{50}$
& mAP$_{50}$
& mAP$_{50}$
& mAP$_{50}$ \\
\hline

YOLO11n
& \textbf{46.0} & 81.9 & 90.4 & 79.7 & 64.7 & 91.5 \\

FFRM-YOLO
& 45.3 & \textbf{85.3} & \textbf{95.6} & \textbf{82.9} & \textbf{67.7} & \textbf{94.8} \\

\hline
\end{tabular*}
\end{table}





\begin{figure}[!t]
    \centering
    \includegraphics[width=0.98\textwidth]{NEU.png}
\caption{Qualitative detection comparison on NEU-DET. Rows show examples labeled by their principal defect category. Columns present ground-truth annotations (GT) and predictions from YOLO11n, MWYOLO, SH-DETR, and FFRM-YOLO. Box colors indicate defect categories, and numbers on predicted boxes denote confidence scores. Some images contain multiple categories.}
    \label{fig:neu_vis}
\end{figure}

To account for signed feature activations, each tensor is converted into a single-channel response map using the root-mean-square across channels, followed by 1st/99th percentile clipping and normalization. The responses are displayed using the Jet colormap with a shared color scale within each row. As shown in Fig.~\ref{fig:subband_vis}, the LL features exhibit broader structural responses in defects such as patches, whereas the high-frequency outputs highlight localized boundaries and texture patterns in scratches, crazing, and rolled-in scale. Since MDRM involves cross-subband interaction, the high-frequency outputs represent refined reconstruction groups rather than isolated directional components. These visualizations illustrate the complementary spatial responses of the two branches, consistent with the detection improvements observed in the ablation studies.

To further examine category-level performance,
Table~\ref{tab:per_class_neu} compares the
per-class mAP$_{50}$ results of YOLO11n and
FFRM-YOLO on NEU-DET. FFRM-YOLO achieves
improvements in five of the six defect categories,
with the largest gain of 5.2 percentage points
observed for Patches. Improvements are also
observed for Inclusion, Pitted Surface, Rolled-in
Scale, and Scratches, while Crazing shows a slight
decrease of 0.7 percentage points. These results
indicate that the proposed frequency-aware
refinement benefits most evaluated defect
categories, although its effectiveness varies
across defect types.


\subsection{Qualitative Analysis}

Fig.~\ref{fig:neu_vis} compares detection results for six defect categories on NEU-DET. The examples were selected from the evaluation set to illustrate representative detection behaviors, including missed detections, overlapping predictions, and multi-label coexistence, rather than through random sampling. All methods were evaluated on the same images and annotations, with selection intended to illustrate detection differences rather than favor a particular method. For scratches and crazing, FFRM-YOLO detects defects missed by YOLO11n, while MWYOLO and SH-DETR occasionally produce overlapping predictions. In the rolled-in scale example, FFRM-YOLO produces predictions more closely aligned with the annotated defects, whereas MWYOLO generates overlapping boxes. The inclusion and patches examples show broadly similar detection patterns across methods, with differences in localization and confidence. The pitted surface example contains annotations for both patches and pitted surface, illustrating multi-label coexistence rather than necessarily indicating false positives. These qualitative observations complement the quantitative results obtained on the full evaluation set.


\section{Discussion}\label{sec:limitations}
Despite its effectiveness, FFRM-YOLO still has limitations in architectural flexibility and model complexity. The framework uses a fixed Haar wavelet basis and predefined convolution kernels, limiting its adaptability to different feature characteristics. Moreover, MDRM exploits directional information from wavelet subbands rather than explicitly employing direction-specific kernels. Incorporating FFRM throughout the backbone and neck also increases computational cost relative to YOLO11n, resulting in lower inference throughput. Although improvements are observed across three datasets, cross-dataset generalization remains unverified, and some architectural ablations require further repeated evaluation. Future work will explore adaptive wavelet representations, efficient feature refinement, and broader generalization evaluations.

\section{Conclusion}
\label{sec:conclusion}
This work proposed FFRM-YOLO, a frequency-domain feature refinement framework for metal surface defect detection based on YOLO11n. Through wavelet decomposition and reconstruction, GLRM and MDRM perform complementary refinement of low-frequency structural information and high-frequency texture details, respectively. Experiments on NEU-DET, GC10-DET, and APDDD achieve $\mathrm{mAP}_{50}$ values of $(78.6 \pm 0.5)\%$, $(66.1 \pm 0.3)\%$, and $(60.6 \pm 0.6)\%$, outperforming YOLO11n by 2.9, 2.7, and 3.0 percentage points in the corresponding mean values. With 4.8~M parameters, 7.7~GFLOPs, and an inference throughput of 135.2~FPS, the model achieves a competitive accuracy--efficiency trade-off. Component ablations and capacity-matched comparisons further support the contributions of frequency-specific refinement and its complementary modules.

\backmatter

\vspace{0.8em}

\noindent\textbf{Author contributions:}\quad


\vspace{0.8em}

\noindent\textbf{Funding:}\quad


\vspace{0.8em}

\noindent\textbf{Data availability:}\quad


\section*{Declarations}

\noindent\textbf{Competing interests:}\quad
The authors declare no competing interests.

% \section*{Declarations}

% Some journals require declarations to be submitted in a standardised format. Please check the Instructions for Authors of the journal to which you are submitting to see if you need to complete this section. If yes, your manuscript must contain the following sections under the heading `Declarations':

% \begin{itemize}
% \item Funding
% \item Conflict of interest/Competing interests (check journal-specific guidelines for which heading to use)
% \item Ethics approval and consent to participate
% \item Consent for publication
% \item Data availability 
% \item Materials availability
% \item Code availability 
% \item Author contribution
% \end{itemize}

% \noindent
% If any of the sections are not relevant to your manuscript, please include the heading and write `Not applicable' for that section. 


% \begin{appendices}

% \section{Section title of first appendix}\label{secA1}

% An appendix contains supplementary information that is not an essential part of the text itself but which may be helpful in providing a more comprehensive understanding of the research problem or it is information that is too cumbersome to be included in the body of the paper.

% \end{appendices}

\begin{thebibliography}{49}

\bibitem{Redmon2016YOLO}
Redmon, J., Divvala, S., Girshick, R. \& Farhadi, A. in \textit{You Only Look Once: Unified, Real-Time Object Detection} 779--788 (IEEE, 2016).

\bibitem{Jocher2024YOLO11}
Jocher, G. \& Qiu, J. \textit{Ultralytics YOLO11} (2024). Software.

\bibitem{Yu2022ESNet}
Yu, X., Lyu, W., Zhou, D., Wang, C. \& Xu, W. ES-Net: Efficient scale-aware network for tiny defect detection. \textit{IEEE Transactions on Instrumentation and Measurement} \textbf{71}, 1--14 (2022).

\bibitem{Liu2025GCNet}
Liu, G., Chu, M., Gong, R. \& Zheng, Z. Global attention module and cascade fusion network for steel surface defect detection. \textit{Pattern Recognition} \textbf{158}, 110979 (2025).

\bibitem{Liu2026GSSAYOLO}
Liu, Y. \& Deng, Y. GSSA-YOLO: A novel lightweight detection method for steel strip surface defect detection. \textit{Measurement Science and Technology} \textbf{37} (2026).

\bibitem{Zheng2026ACMSPYOLO}
Zheng, L. et al. A lightweight steel surface defect detection model for automotive manufacturing production lines. \textit{Engineering Applications of Artificial Intelligence} \textbf{175}, 114601 (2026).

\bibitem{Zhang2026DRFMFF}
Zhang, H. et al. Steel surface defect detection based on dynamic receptive field
and multi-scale features fusion. Sci. Rep. 10.1038/s41598-026-70164-7 (2026).

\bibitem{Zhang2024MSFCRA}
Zhang, H. et al. Surface defect detection of hot rolled steel based on multi-scale feature fusion and attention mechanism residual block. \textit{Scientific Reports} \textbf{14}, 7671 (2024).

\bibitem{Mallat1989Wavelet}
Mallat, S. G. A theory for multiresolution signal decomposition: The wavelet representation. \textit{IEEE Transactions on Pattern Analysis and Machine Intelligence} \textbf{11}, 674--693 (1989).

\bibitem{Liu2018MWCNN}
Liu, P., Zhang, H., Zhang, K., Lin, L. \& Zuo, W. in \textit{Multi-Level Wavelet-CNN for Image Restoration} 773--782 (IEEE, 2018).

\bibitem{Cui2025SWT}
Cui, Q., Li, Y., Bian, H., Kong, J. \& Dong, Y. Visual defect recognition with stationary wavelet transform based neural networks. \textit{Digital Signal Processing} \textbf{158}, 104947 (2025).

\bibitem{Chen2026WTCFNet}
Chen, H. \& Ren, Y.-B. WTCFNet for industrial defect detection using wavelet transform and cross layer feature fusion. \textit{Scientific Reports} \textbf{16}, 23579 (2026).

\bibitem{Gu2024Mamba}
Gu, A. \& Dao, T. in \textit{Mamba: Linear-Time Sequence Modeling with Selective State Spaces} (OpenReview.net, 2024).

\bibitem{Liu2024VMamba}
Liu, Y. et al. in \textit{VMamba: Visual State Space Model}, Vol. 37, 103031--103063 (Curran Associates, Inc., 2024).

\bibitem{Zhang2025WaveMambaAD}
Zhang, Q., Shao, M., Chen, X., Lv, X. \& Xu, K. in \textit{Wave-MambaAD: Wavelet-Driven State Space Model for Multi-Class Unsupervised Anomaly Detection} 20868--20877 (IEEE, 2025).

\bibitem{Chen2026MWYOLO}
Chen, J., Ren, Z., Xiao, H. \& Wu, Z. MWYOLO: A Mamba-enhanced lightweight YOLO framework with multi-frequency attention for industrial surface defect detection. \textit{AI Engineering} \textbf{1}, 3 (2026).

\bibitem{Ge2026WaveMambaYOLO}
Ge, A., Lv, Y. \& Huang, J. WaveMamba-YOLO: Combining frequency awareness and state-space modeling for defect localization. \textit{PLOS ONE} \textbf{21}, e0344940 (2026).

\bibitem{Zhang2023InnerIoU}
Zhang, H., Xu, C. \& Zhang, S. Inner-IoU: More effective intersection over union loss with auxiliary bounding box. \textit{arXiv preprint arXiv:2311.02877} (2023).

\bibitem{Cai2018CascadeRCNN}
Cai, Z. \& Vasconcelos, N. in \textit{Cascade R-CNN: Delving Into High Quality Object Detection} 6154--6162 (IEEE, 2018).

\bibitem{He2017MaskRCNN}
He, K., Gkioxari, G., Doll{\'a}r, P. \& Girshick, R. in \textit{Mask R-CNN} 2961--2969 (IEEE, 2017).

\bibitem{Carion2020DETR}
Carion, N. et al. in \textit{End-to-End Object Detection with Transformers} 213--229 (Springer, 2020).

\bibitem{Wu2025SHDETR}
Wu, S. et al. SH-DETR: Enhancing steel surface defect detection and classification with an improved transformer architecture. \textit{PLOS ONE} \textbf{20}, e0334048 (2025).

\bibitem{Zhou2025MESCDETR}
Zhou, S., Cai, Y., Zhang, Z. \& Yin, J. MESC-DETR: An improved RT-DETR algorithm for steel surface defect detection. \textit{Electronics} \textbf{14}, 2232 (2025).

\bibitem{Hsiang2025RTLSMDETR}
Hsiang, H.-C. \& Zheng, X. in \textit{Real-Time Steel Surface Defect Detection Using RT-LSM-DETR}, Vol. 13681, 1368129 (SPIE, 2025).

\bibitem{Yu2026HMADETR}
Yu, J., Chen, Y., Ma, F. \& Sun, Y. HMA-DETR: A real-time hybrid multi-scale adaptive feature learning network for steel surface defect detection. \textit{Journal of Real-Time Image Processing} \textbf{23}, 150 (2026).

\bibitem{Luo2025LiteDETR}
Luo, T. et al. A lightweight defect detection transformer for printed circuit boards combining image feature augmentation and refined cross-scale feature fusion. \textit{Engineering Applications of Artificial Intelligence} \textbf{155}, 111128 (2025).

\bibitem{Li2024EWT}
Li, J. et al. EWT: Efficient wavelet-transformer for single image denoising. \textit{Neural Networks} \textbf{177}, 106378 (2024).

\bibitem{Zou2024WaveMamba}
Zou, W., Gao, H., Yang, W. \& Liu, T. in \textit{Wave-Mamba: Wavelet State Space Model for Ultra-High-Definition Low-Light Image Enhancement} (Association for Computing Machinery, 2024).

\bibitem{Gardner2026WEPaDiM}
Gardner, C., Dhiren, A., Min, B. \& Ahn, T.-H. Wavelet-enhanced PaDiM for industrial anomaly detection. \textit{IEEE Access} \textbf{14}, 37702--37718 (2026).

\bibitem{Yan2026WFSCAYOLO}
Yan, J. et al. WFSCA-YOLO: Robust object detection for terrestrial optical sensing under atmospheric degradation via a wavelet-driven frequency--spatial co-awareness framework. \textit{Remote Sensing} \textbf{18}, 1667 (2026).

\bibitem{Huang2024LocalMamba}
Huang, T. et al. in \textit{LocalMamba: Visual State Space Model with Windowed Selective Scan} 12--22 (Springer, 2024).

\bibitem{He2024MambaAD}
He, H. et al. in \textit{MambaAD: Exploring State Space Models for Multi-Class Unsupervised Anomaly Detection}, Vol. 37 (Curran Associates, Inc., 2024).

\bibitem{Lei2025YOLOGVMamba}
Lei, C., Wang, C., Xing, Z. \& Bi, X. in \textit{YOLO-GVMamba: An Efficient Steel Surface Defect Object Detection Method Based on State Space Model} (IEEE, 2025).

\bibitem{He2025MobileMamba}
He, H. et al. in \textit{MobileMamba: Lightweight Multi-Receptive Visual Mamba Network} 4497--4507 (IEEE, 2025).

\bibitem{Xu2026SCSMamba}
Xu, Y., Zhao, X. \& Xiao, X. SCS-Mamba: A lightweight Sobel-enhanced content-aware state space model with shared deformable adaptive head for efficient underwater object detection. \textit{Journal of Real-Time Image Processing} \textbf{23}, 4 (2026).

\bibitem{Fu2026LVMSOD}
Fu, Y. et al. LVMSOD: Lightweight visual Mamba small object detection for autonomous vehicles. \textit{IEEE Internet of Things Journal} \textbf{13}, 19649--19666 (2026).

\bibitem{Hu2018SENet}
Hu, J., Shen, L. \& Sun, G. in \textit{Squeeze-and-Excitation Networks} 7132--7141 (IEEE, 2018).

\bibitem{Zheng2020distance}
Zheng, Z. et al. in \textit{Distance-IoU Loss: Faster and Better Learning for Bounding Box Regression}, Vol. 34, 12993--13000 (AAAI Press, 2020).

\bibitem{Li2020GFL}
Li, X. et al. in \textit{Generalized Focal Loss: Learning Qualified and Distributed Bounding Boxes for Dense Object Detection}, Vol. 33 (Curran Associates, Inc., 2020).

\bibitem{Song2013NEU}
Song, K. \& Yan, Y. A noise robust method based on completed local binary patterns for hot-rolled steel strip surface defects. \textit{Applied Surface Science} \textbf{285}, 858--864 (2013).

\bibitem{Lv2020GC10}
Lv, X., Duan, F., Jiang, J.-J., Fu, X. \& Gan, L. Deep metallic surface defect detection: The new benchmark and detection network. \textit{Sensors} \textbf{20}, 1562 (2020).

\bibitem{tianchi2018aluminium}
Tianchi, Alibaba Cloud. Guangdong industrial intelligence big data innovation competition: Aluminium profile surface defect detection dataset. Available at: https://tianchi.aliyun.com/competition/entrance/231682/information (2018).

\bibitem{Loshchilov2019AdamW}
Loshchilov, I. \& Hutter, F. in \textit{Decoupled Weight Decay Regularization} (OpenReview.net, 2019).

\bibitem{Jocher2020YOLOv5}
Jocher, G. \textit{Ultralytics YOLOv5} (2020). Software.

\bibitem{Jocher2023YOLOv8}
Jocher, G., Chaurasia, A. \& Qiu, J. \textit{Ultralytics YOLOv8} (2023). Software.

\bibitem{Wang2024YOLOv9}
Wang, C.-Y., Yeh, I.-H. \& Liao, H.-Y. M. in \textit{YOLOv9: Learning What You Want to Learn Using Programmable Gradient Information} 1--21 (Springer, 2024).

\bibitem{Wang2024YOLOv10}
Wang, A. et al. in \textit{YOLOv10: Real-Time End-to-End Object Detection}, Vol. 37 (Curran Associates, Inc., 2024).

\bibitem{Zheng2024MDYOLO}
Zheng, H., Chen, X., Cheng, H., Du, Y. \& Jiang, Z. MD-YOLO: Surface defect detector for industrial complex environments. \textit{Optics and Lasers in Engineering} \textbf{178}, 108170 (2024).

\bibitem{Deng2026SDCMYOLO}
Deng, F., Chen, W. \& Fu, L. SDCM-YOLO: A cross-scale, texture-aware framework for metal surface defect detection. \textit{Engineering Research Express} \textbf{8}, 075211 (2026).

\end{thebibliography}

\end{document}
